\documentclass[conference]{IEEEtran}
\IEEEoverridecommandlockouts
\usepackage{cite}
\usepackage{amsmath,amssymb,amsfonts}
\usepackage{array}
\usepackage{tabularx}
\usepackage{textcomp}
\usepackage{xcolor}
\usepackage{graphicx}
\usepackage{microtype}
\def\BibTeX{{\rm B\kern-.05em{\sc i\kern-.025em b}\kern-.08em
    T\kern-.1667em\lower.7ex\hbox{E}\kern-.125emX}}

\begin{document}

\title{Energy-Efficient Attention Accelerators for Transformer Inference: \\
A Survey of Software-Hardware Co-Design}

\author{\IEEEauthorblockN{Author Name}
\IEEEauthorblockA{\textit{Department, Institution} \\
City, Country \\
email@institution.edu}
}

\maketitle

\begin{abstract}
Transformer-based large language model (LLM) inference is increasingly limited by attention and key-value (KV) cache data movement rather than peak arithmetic throughput, especially under long contexts and token-by-token decode.
This survey reviews monolithic accelerator research for Transformer inference---covering both prompt prefill and decode---through four complementary software--hardware themes: FlashAttention-native dataflows and array architectures; mixed-precision attention datapaths and KV-cache quantization; specialized hardware for non-general-matrix-multiply operators such as softmax, root-mean-square normalization, and rotary position embedding; and compilation, mapping, and evaluation toolchains.
We summarize how IO-aware online softmax, low-width KV storage with higher-precision accumulation, approximate normalization datapaths, and phase-aware mappers recur across recent GPU kernels and application-specific integrated circuit designs, and we compare proposals along unified metrics including processing-element utilization, high-bandwidth memory traffic, energy per token, and accuracy.
Although each theme has progressed rapidly, most published systems optimize only one or two concerns on mismatched platforms.
We identify integrated single-chip co-design---jointly validating dataflow, mixed precision, non-general-matrix-multiply fusion, and compile-time scheduling through algorithmic, architectural, and register-transfer-level evaluation---as the central open problem for energy-efficient attention acceleration, and we briefly situate sparse and memory-centric scale-out approaches as adjacent comparators outside this taxonomy.
\end{abstract}

\begin{IEEEkeywords}
Attention accelerator, FlashAttention, KV cache, mixed precision, systolic array, software--hardware co-design, LLM inference.
\end{IEEEkeywords}

\section{Introduction and Background}
\label{sec:intro}

Transformer-based large language models (LLMs)~\cite{vaswani2017attention} now dominate large-scale natural-language inference.
Deployers increasingly optimize end-to-end latency, throughput, and energy per generated token rather than peak floating-point throughput alone.
In each Transformer layer, multi-head self-attention (MHA) together with the key-value (KV) cache forms the critical path for both \emph{prefill}, which processes the prompt in parallel, and \emph{decode}, which emits one token at a time.
Attention acceleration is therefore a distinct software--hardware co-design problem---not reducible to accelerating generic general matrix multiply (GEMM) on systolic or tensor engines~\cite{chen2016eyeriss,ma2024tpuv4i}.

\subsection{Motivation and the Attention Bottleneck}
\label{sec:intro-motivation}

The dominant bottleneck in LLM inference has shifted from arithmetic throughput to data movement.
For head $h$ with dimension $d_h$, MHA computes
\begin{align}
Q_h &= X W^Q_h,\quad K_h = X W^K_h,\quad V_h = X W^V_h, \nonumber\\
S_h &= \frac{Q_h K_h^\top}{\sqrt{d_h}},\quad
O_h = \mathrm{softmax}(S_h)\,V_h,
\label{eq:mha}
\end{align}
where $X \in \mathbb{R}^{N \times d_{\mathrm{model}}}$ is the layer input and $W^Q_h, W^K_h, W^V_h$ are learned projections~\cite{vaswani2017attention}.
During decode, only the final row of $Q_h$ is active while $K_h$ and $V_h$ are read from a growing KV cache of length $N$; during prefill, $Q_h$, $K_h$, and $V_h$ share prompt length $N_p$ and the cache is populated once.
Prefill is more sensitive to $\mathcal{O}(N_p^2 d_h)$ attention work, whereas decode becomes increasingly \emph{memory-bound} as $\mathcal{O}(N d_h)$ cached data must be streamed from high-bandwidth memory (HBM) each step~\cite{dao2022flashattention,dao2023flashdecoding,bitdecoding2026}.

Naively materializing $S_h \in \mathbb{R}^{N \times N}$ in HBM incurs prohibitive traffic at long context lengths.
FlashAttention and follow-on GPU kernels mitigate this cost through IO-aware tiling and online softmax, preserving numerical equivalence to~\eqref{eq:mha} while reducing off-chip accesses~\cite{dao2022flashattention,dao2023flashattention2,shah2024flashattention3}.
Dedicated accelerators must still realize the same dataflow under fixed on-chip static random-access memory (SRAM) and processing-element (PE) budgets, fuse non-GEMM stages that otherwise stall systolic pipelines~\cite{wang2023cosa,desa2025}, and compress KV cache traffic through mixed-precision datapaths~\cite{liu2024kivi,sawint42026,bitdecoding2026}.

\subsection{FlashAttention and Quantization Baselines}
\label{sec:intro-baselines}

\paragraph{IO-aware attention dataflow}
FlashAttention~\cite{dao2022flashattention} processes $K_h$ and $V_h$ in tiles while maintaining per-row statistics for \emph{online softmax}.
For query row $i$, with running maximum $m_i$ and exponential sum $\ell_i$ after blocks $1,\ldots,j{-}1$,
\begin{align}
m_i' &= \max(m_i, \max_k S_{ij,k}), \label{eq:online-max}\\
\ell_i' &= e^{m_i - m_i'} \ell_i + \textstyle\sum_k e^{S_{ij,k} - m_i'}, \label{eq:online-sum}\\
O_i' &= e^{m_i - m_i'} O_i + \textstyle\sum_k e^{S_{ij,k} - m_i'} V_{j,k}, \label{eq:online-o}
\end{align}
where $O_i$ is the partial output accumulator.
Subtracting the row maximum before exponentiation keeps all exponent inputs $\leq 0$, stabilizing softmax and easing hardware $\exp(\cdot)$ approximation~\cite{lin2025systolicattention,koca2023exp}.
FlashAttention-2 improves GPU parallelism~\cite{dao2023flashattention2}; FlashAttention-3 adds asynchrony and FP8 on Hopper-class GPUs~\cite{shah2024flashattention3,micikevicius2022fp8}.
FlashDecoding~\cite{dao2023flashdecoding} applies the same IO-aware principle to decode via KV splitting.
These kernels serve as \emph{algorithmic gold references} for correctness, traffic complexity, and GPU throughput.

\paragraph{Non-GEMM operators}
Attention blocks also require row-wise softmax (max, $\exp$, reciprocal), root-mean-square normalization (RMSNorm)~\cite{zhang2019root}, and rotary position embedding (RoPE)~\cite{su2024roformer} on $Q$ and $K$.
Although individually smaller than GEMM, they frequently break PE utilization when mapped to external vector units~\cite{wang2023cosa,mive2026,sole2025}.

\paragraph{Mixed-precision and KV compression}
Low-bit formats reduce KV-cache footprint but can distort $QK^\top$ products and softmax accumulation; examples include FP8~\cite{micikevicius2022fp8}, INT4, and microscaling MXFP4.
Widely used post-training quantization (PTQ) baselines are SmoothQuant~\cite{xiao2023smoothquant}, GPTQ~\cite{frantar2023gptq}, QuaRot~\cite{ashkboos2024quarot}, and SpinQuant~\cite{liu2024spinquant}.
KV-oriented schemes include KIVI~\cite{liu2024kivi}, SAW-INT4~\cite{sawint42026}, and BitDecoding~\cite{bitdecoding2026}.
Section~\ref{sec:l2} develops the recurring rule of storing $K/V$ at low width while keeping matmul and softmax at higher precision and fusing dequantization with MAC.

Table~\ref{tab:baselines} summarizes representative baselines referenced throughout the survey~\cite{dao2022flashattention,dao2023flashattention2,shah2024flashattention3,dao2023flashdecoding,lin2025systolicattention,plena2025,bitdecoding2026,sawint42026,wang2023cosa}.
Architectural details and cross-design comparisons appear in Sections~\ref{sec:l1}--\ref{sec:l4}; absolute performance numbers are not directly comparable across GPU and ASIC platforms.

\begin{table}[t]
\caption{Representative attention acceleration baselines.
L1--L4 map to Sections~\ref{sec:l1}--\ref{sec:l4}.}
\label{tab:baselines}
\centering
\footnotesize
\setlength{\tabcolsep}{4pt}
\begin{tabularx}{\columnwidth}{@{}l c >{\raggedright\arraybackslash}X c@{}}
\hline
\textbf{Work} & \textbf{Plat.} & \textbf{Role} & \textbf{Thm.} \\
\hline
FA v1--v3 & GPU & IO-aware dataflow reference & L1 \\
FlashDecoding & GPU & Decode-phase throughput baseline & L1 \\
SystolicAttention & ASIC & Native FlashAttention in one array & L1 \\
PLENA & ASIC & Flat array with mixed precision and compiler & L1--L4 \\
BitDecoding & GPU & Low-bit KV cache for decode & L2 \\
SAW-INT4 & GPU & Block rotation with INT4 KV cache & L2 \\
COSA+ & ASIC & Dual cooperative systolic arrays with softmax unit & L1 \\
\hline
\end{tabularx}
\end{table}

\subsection{Survey Scope, Taxonomy, and Gap}
\label{sec:intro-scope}

This survey covers Transformer \emph{inference} on single-chip or monolithic accelerators, including both prefill and decode, and includes works with explicit hardware organizations or algorithm--hardware co-design for the attention datapath.
We do not treat distributed training, multi-chiplet scale-out, processing-in-memory (PIM) or compute-in-memory (CIM) centric designs, or sparse mixture-of-experts (MoE) scheduling as primary topics; Section~\ref{sec:adjacent} cites representative adjacent works only to clarify boundaries~\cite{salca2026,liu2021sanger,amma2026,lolpim2024,heo2024nuepim}.

We organize the literature around four complementary themes:
\emph{(L1) FlashAttention-native dataflows and array architectures} map online softmax and tiled $QK^\top$/softmax/$AV$ pipelines onto systolic or flattened arrays~\cite{lin2025systolicattention,plena2025,flatattention2026,wang2023cosa,streamattention2025,desa2025};
\emph{(L2) mixed-precision attention datapaths and KV cache quantization} combine FP8, INT4, or MXFP4 with rotation-based outlier suppression and fused dequantization~\cite{xiao2023smoothquant,ashkboos2024quarot,liu2024kivi,sawint42026,bitdecoding2026,batquant2026,ultraquant2026};
\emph{(L3) non-GEMM specialized units} implement softmax, RMSNorm, and RoPE with dedicated or approximate hardware~\cite{mive2026,sole2025,sun2022softmax,koca2023exp};
\emph{(L4) compilation, mapping, and evaluation toolchains} automate tiling, direct memory access (DMA) overlap, and PPA or register-transfer-level (RTL) assessment~\cite{parashar2019timeloop,wu2019accelergy,scalesimv3_2025,klhufek2025transinfersim,tilelang2025,plena2025}.

Although each theme has advanced quickly, synthesis remains difficult.
GPU kernels, single-array FlashAttention engines, KV quantization schemes, and standalone non-GEMM accelerators often optimize one or two concerns in isolation~\cite{lin2025systolicattention,liu2024kivi,mive2026,plena2025}.
Cross-study comparison is further complicated by mismatched benchmarks and ambiguous separation of prefill versus decode results.
Less frequently addressed is unified co-optimization of dataflow mapping, mixed precision, non-GEMM fusion, and compile-time scheduling on one monolithic accelerator; Section~\ref{sec:gaps} examines this gap.

\subsection{Comparison Framework and Paper Organization}
\label{sec:intro-metrics}

Across Sections~\ref{sec:l1}--\ref{sec:l4}, we characterize designs along shared dimensions: problem entry point, dataflow and precision scheme, hardware structure, evaluation level (algorithmic, architectural, or RTL), reported latency, PE utilization, HBM traffic, energy per token, accuracy, open-source availability, and relevance to integrated co-design.
Table~\ref{tab:metrics} lists unified metrics; we separate prefill and decode whenever source data permit~\cite{dao2023flashdecoding,bitdecoding2026}.
Energy and throughput gains are interpreted as \emph{relative} improvements at iso-accuracy unless authors report absolute silicon data.
Accuracy--efficiency trade-offs for low-bit KV schemes are best read as Pareto fronts~\cite{sawint42026,batquant2026}.

\begin{table}[t]
\caption{Unified metrics for cross-study comparison.}
\label{tab:metrics}
\centering
\small
\begin{tabular}{@{}l l@{}}
\hline
\textbf{Category} & \textbf{Representative metrics} \\
\hline
Performance & Latency; throughput (tokens/s) \\
Utilization & PE utilization (effective / peak MAC) \\
Memory & SRAM traffic; HBM traffic (bytes) \\
Energy \& area & Energy/token; TOPS/W; area (mm$^2$) \\
Accuracy & Task accuracy; PPL; cosine sim.\ vs.\ FP32 \\
\hline
\end{tabular}
\end{table}

The remainder of this paper is organized as follows.
Sections~\ref{sec:l1}--\ref{sec:l4} review the four themes above.
Section~\ref{sec:adjacent} discusses sparse and memory-centric adjacent directions.
Section~\ref{sec:gaps} synthesizes research gaps and outlook.
Section~\ref{sec:conclusion} concludes.

\section{FlashAttention-Native Dataflow and Array Architectures}
\label{sec:l1}

The first survey theme (L1) examines how \emph{IO-aware} attention dataflows---especially the online softmax of FlashAttention~\cite{dao2022flashattention,dao2023flashattention2,shah2024flashattention3}---are mapped onto fixed on-chip arrays rather than GPU shared-memory hierarchies.
We include monolithic ASIC and tile-based accelerators with explicit attention datapath co-design, and treat GPU FlashAttention and FlashDecoding kernels as algorithmic and throughput references (Section~\ref{sec:intro-baselines}) rather than hardware organizations.
Sparse or dynamic attention variants (e.g., SpAtten~\cite{wang2021spatten}) are discussed only in Section~\ref{sec:adjacent}.

\subsection{Online Softmax as a Hardware Primitive}
\label{sec:l1-online}

Section~\ref{sec:intro-baselines} introduced the online update rules in~\eqref{eq:online-max}--\eqref{eq:online-o}.
For dedicated hardware, three properties matter beyond numerical equivalence to~\eqref{eq:mha}.
First, row-wise statistics ($m_i$, $\ell_i$, partial $O_i$) can be updated \emph{in place} as $K/V$ tiles stream through the array, avoiding materialization of $S_h \in \mathbb{R}^{N \times N}$ in high-bandwidth memory (HBM).
Second, subtracting the running row maximum before exponentiation bounds all $\exp$ inputs at $\leq 0$, which stabilizes softmax and enables piecewise-linear or split-integer $\exp$ approximations that reuse multiply-accumulate (MAC) datapaths~\cite{lin2025systolicattention,koca2023exp}.
Third, the same reduction pattern---compare-and-update max, accumulate weighted sums---maps naturally onto systolic upward propagation or comparator trees, whereas naively offloading softmax to an external vector unit typically stalls the main GEMM pipeline~\cite{wang2023cosa,desa2025,streamattention2025}.

\subsection{Hardware Mapping Taxonomy}
\label{sec:l1-taxonomy}

Recent accelerators cluster into five mapping patterns along the $QK^\top \rightarrow \mathrm{softmax} \rightarrow AV$ pipeline; buffer sizes and tile orders differ by design, but the organizational split is stable across recent proposals~\cite{wang2023cosa,lin2025systolicattention,plena2025,flatattention2026,streamattention2025,desa2025}.

\paragraph{GPU IO-aware kernels}
FlashAttention v1--v3 and FlashDecoding~\cite{dao2022flashattention,dao2023flashattention2,shah2024flashattention3,dao2023flashdecoding} tile $K/V$ through shared memory, overlap memory traffic with tensor-core GEMM, and expose the reference traffic complexity and decode-phase KV splitting for long contexts.
GPU-side operator fusion such as FLAT~\cite{kao2023flat} optimizes attention scheduling in software before custom silicon is engaged.
They are the correctness and GPU-throughput baselines against which ASIC dataflows are judged, but their mapping assumptions (warp-level parallelism, deep software-managed caches) do not transfer directly to shallow on-chip SRAM budgets.

\paragraph{Dual-array with dedicated softmax}
COSA+~\cite{wang2023cosa} assigns separate cooperating systolic arrays to $QK^\top$ and $AV$, with a dedicated softmax unit between them.
This decouples GEMM tiling from softmax latency and reportedly achieves high multi-head attention (MHA) layer utilization ($\sim$95\% in architectural simulation), but the intermediate score buffering and unit-to-unit handoff add area and scheduling complexity relative to fused single-array designs.

\paragraph{Single systolic array with fused reductions}
SystolicAttention (FSA)~\cite{lin2025systolicattention} and StreamAttention~\cite{streamattention2025} keep the full FlashAttention loop inside one array.
FSA augments a $128 \times 128$ systolic mesh with upward data paths and comparator arrays for row-max/row-sum, plus a split-integer $\exp$ unit that reuses MAC hardware; authors report $\sim$4.83$\times$ and $\sim$1.77$\times$ attention throughput versus TPUv5e and Neuron-v2, respectively, at $\sim$12\% area overhead from 16\,nm RTL synthesis.
StreamAttention pursues the same goal through a four-stage streaming operand schedule and Chebyshev-polynomial $\exp$, targeting higher sustained utilization than prior single-array FlashAttention engines that reportedly stall at $\sim$40\% PE occupancy.

\paragraph{Flattened array with head-level decomposition}
PLENA~\cite{plena2025} replaces a rigid 2-D systolic grid with a \emph{flattened} array and head-level decomposition to balance prefill and decode.
It natively executes FlashAttention while co-designing asymmetric quantization pathways (developed further in Section~\ref{sec:l2}) and a custom instruction-set architecture (ISA) with compiler support (Section~\ref{sec:l4}).
Reported end-to-end gains include $\sim$2.23$\times$ throughput versus A100 and $\sim$4.70$\times$ versus TPU v6e, with $\sim$4.04$\times$ energy efficiency over A100 at iso-accuracy in the authors' agentic long-context benchmarks.

\paragraph{Tile fabric with collective optimization}
FlatAttention~\cite{flatattention2026} targets tile-based and wafer-scale organizations where on-chip network \emph{fabric collectives} must be co-optimized with attention tiling.
By aligning data movement with multicast/reduction patterns across $32 \times 32$ tiles, the design reports up to $\sim$92.3\% utilization, $\sim$4.1$\times$ speedup versus FlashAttention-3, and up to $\sim$16$\times$ HBM traffic reduction relative to GPU baselines in architectural simulation---illustrating that dataflow choices at the network level can dominate raw MAC throughput for long-context inference.

\paragraph{Score-stationary hybrids}
DESA~\cite{desa2025} couples a score-stationary systolic array with an immediate-processing unit (IPU) and reconfigurable vector logic for non-GEMM stages that would otherwise idle the mesh during Transformer inference.
Relative to COSA+'s dedicated softmax engine and FSA's in-array reductions, DESA prioritizes layer-wide \emph{dataflow efficiency}.
It improves PE utilization over GPU baselines, yet still forwards part of the non-GEMM work to auxiliary vector paths (Section~\ref{sec:l3}).

Table~\ref{tab:l1} summarizes representative L1 designs along the unified dimensions from Section~\ref{sec:intro-metrics}.
Absolute latency and energy numbers are not directly comparable across unlike platforms; we emphasize architectural choices and \emph{relative} claims at iso-accuracy where authors report them.

\begin{table}[t]
\caption{Representative FlashAttention-native accelerator designs (L1).
Util.\ = PE or layer utilization; eval.\ = dominant evaluation level.}
\label{tab:l1}
\centering
\scriptsize
\setlength{\tabcolsep}{2pt}
\renewcommand{\arraystretch}{1.05}
\resizebox{\columnwidth}{!}{%
\begin{tabular}{@{}lllll@{}}
\hline
\textbf{Work} & \textbf{Array org.} & \textbf{Softmax} & \textbf{Util.} & \textbf{Eval.} \\
\hline
FA/FlashDec. & GPU shared-mem tile & Online exact & {---} & GPU \\
COSA+ & Dual systolic & Dedicated & {$\sim$95\%} & Arch. \\
SystolicAttn & $128^2$ systolic & In-array exp & High & RTL \\
StreamAttn & 4-stage stream & PWL in-array & {$\gg$40\%} & Arch. \\
PLENA & Flat array & Native FA & Bal.~P/D & A{+}RTL \\
FlatAttn & $32^2$ tile fabric & Tiled online & {$\sim$92\%} & Arch. \\
DESA & Score-stat.{+}IPU & IPU-assisted & Impr. & Arch. \\
\hline
\end{tabular}%
}
\end{table}

\subsection{Prefill versus Decode Mapping}
\label{sec:l1-prefill-decode}

Prefill and decode impose different bottlenecks on the same mathematical attention, and L1 designs address them unevenly.
During \emph{prefill}, $Q$, $K$, and $V$ share prompt length $N_p$; attention work scales as $\mathcal{O}(N_p^2 d_h)$ per head and is often compute- or SRAM-capacity limited if tiles fit on chip~\cite{dao2022flashattention,lin2025systolicattention}.
FlashAttention-style tiling reduces HBM traffic proportional to $N_p^2$ score materialization, and wide systolic arrays favor batched prompt processing.

During \emph{decode}, only one new query row is active per step while the KV cache grows linearly with context length $N$.
Each token requires streaming $\mathcal{O}(N d_h)$ cached $K/V$ from HBM, making decode increasingly memory-bound~\cite{dao2023flashdecoding,bitdecoding2026}.
FlashDecoding mitigates this on GPUs by parallelizing over KV splits; PLENA explicitly rebalances head-level mapping between phases~\cite{plena2025}; FlatAttention reduces cross-tile HBM movement via fabric collectives~\cite{flatattention2026}.
Few monolithic ASIC proposals publish phase-separated latency or bandwidth numbers, which complicates cross-study comparison and motivates reporting prefill and decode metrics independently (Table~\ref{tab:metrics}).

\subsection{Open Problems}
\label{sec:l1-open}

Three gaps recur across L1 literature.
\emph{(i) Phase-aware unified mapping.} Designs that excel at prefill tiling (wide $QK^\top$ blocks) may under-utilize the array during skinny decode matmuls unless the mapper adapts tile shapes, array folding, or head batching~\cite{plena2025,streamattention2025,dao2023flashdecoding}.
\emph{(ii) Non-GEMM fusion depth.} In-array softmax (FSA, StreamAttention) reduces external vector stalls, but RMSNorm and rotary position embedding (RoPE) remain largely outside L1 scope and are treated in Section~\ref{sec:l3}.
\emph{(iii) Precision co-design.} Most L1 ASIC studies assume FP16/FP32 attention internals; coupling native FlashAttention tiling with low-bit KV datapaths (Section~\ref{sec:l2}) without breaking online softmax invariants is still largely unexplored on monolithic silicon beyond PLENA's initial pathways~\cite{plena2025,liu2024kivi,sawint42026}.
Joint optimization of L1 dataflow, L2 mixed precision, L3 non-GEMM units, and L4 compile-time scheduling on one accelerator remains the central opportunity discussed in Section~\ref{sec:gaps}.

\section{Mixed-Precision Attention Datapath and KV Cache Quantization}
\label{sec:l2}

The second survey theme (L2) addresses how attention accelerators reduce KV-cache footprint and HBM traffic through mixed-precision datapaths while preserving acceptable end-to-end accuracy.
Section~\ref{sec:intro-baselines} previewed the recurring design rule: \emph{store $K$ and $V$ at low width, execute $QK^\top$, online softmax, and $AV$ at higher precision, and fuse dequantization with MAC wherever possible}~\cite{liu2024kivi,sawint42026,bitdecoding2026,plena2025}.
Most published L2 work targets GPU serving kernels or algorithmic post-training quantization (PTQ); monolithic ASIC datapaths with native low-bit KV remain comparatively rare and are discussed jointly with PLENA in Sections~\ref{sec:l1} and~\ref{sec:l4}.

\subsection{Precision Formats and Error Budgets}
\label{sec:l2-formats}

Attention quantization differs from weight-only compression because errors propagate through three coupled stages: the $QK^\top$ dot products that form logits, the row-wise softmax that amplifies relative score errors, and the $AV$ reduction that accumulates value errors across context length~\cite{dao2022flashattention,micikevicius2022fp8,bitdecoding2026}.
Formats in recent L2 literature span FP16 baselines, FP8 tensor operands~\cite{micikevicius2022fp8,shah2024flashattention3}, INT8 weight--activation PTQ~\cite{xiao2023smoothquant}, INT4 (and sub-4-bit) KV caches~\cite{liu2024kivi,sawint42026}, and microscaling block formats such as MXFP4/NVFP4~\cite{bitdecoding2026,batquant2026}.
Block scales---per-channel, per-token, or fixed-size microscaling groups (e.g., 32 elements)---trade metadata overhead against outlier resilience; misaligned scale granularity relative to rotation or smoothing transforms is a recurring source of accuracy loss~\cite{ashkboos2024quarot,batquant2026,ultraquant2026}.

For hardware mapping, three numerical constraints recur.
First, softmax stability still requires row-max tracking even when logits are computed from low-bit $K$ or dequantized on the fly~\cite{lin2025systolicattention,liu2024kivi}.
Second, asymmetric treatment of $K$ versus $V$ is common: KIVI quantizes keys per channel and values per token at 2-bit width~\cite{liu2024kivi}, reflecting unequal outlier structure in cached tensors.
Third, decode-phase sensitivity dominates serving evaluations because each generated token re-reads the full compressed cache, so bandwidth savings and dequantization latency must be reported separately from prefill~\cite{bitdecoding2026,sawint42026,ultraquant2026,dao2023flashdecoding}.

\subsection{Quantization Method Taxonomy}
\label{sec:l2-taxonomy}

We group L2 methods by \emph{what} they quantize and \emph{how} they suppress outliers, rather than by publication venue.

\paragraph{Weight--activation PTQ baselines}
SmoothQuant~\cite{xiao2023smoothquant} migrates activation difficulty into weights via per-channel smoothing, enabling W8A8 inference with near-FP16 accuracy on language models.
GPTQ~\cite{frantar2023gptq} compresses weights post hoc using Hessian-aware rounding.
These schemes establish INT8/FP8 baselines for linear layers but do not by themselves solve KV-cache growth; they appear in L2 surveys mainly as precursors to aggressive KV formats and as accuracy anchors in combined W4A4/KV4 evaluations~\cite{batquant2026}.

\paragraph{Rotation-based outlier suppression}
QuaRot~\cite{ashkboos2024quarot} applies a global Hadamard rotation so that INT4 weights, activations, and KV tensors encounter fewer extreme values; rotations can be fused into GEMM datapaths.
SpinQuant~\cite{liu2024spinquant} learns rotation matrices for tighter INT4/INT8 grids.
Global rotations align well with uniform INT4 grids but can mismatch microscaled MXFP4 blocks; BATQuant~\cite{batquant2026} addresses this with block-wise affine transforms (BAT) aligned to MXFP granularity, reporting improved W4A4 and W4A8/KV8 accuracy over round-to-nearest and QuaRot-style global rotation on lm-eval benchmarks.

\paragraph{KV-cache-specific schemes}
KIVI~\cite{liu2024kivi} introduces tuning-free asymmetric 2-bit KV compression with higher-precision residual handling, reducing memory at moderate accuracy cost.
SAW-INT4~\cite{sawint42026} targets \emph{real serving} constraints with token-wise INT4 KV caches and block-diagonal Hadamard rotation (BDR) to curb outliers without global rotation overhead; authors integrate prefill with FlashAttention-3 and decode with fused Triton kernels while preserving paged-KV layouts, reporting task accuracy on GSM8K, GPQA, and MATH500 comparable to plain INT4 at similar throughput.
UltraQuant~\cite{ultraquant2026} focuses on multi-turn agent workloads, combining FP4 KV, FP8 queries, UE8M0 group scales, and TurboQuant-style rotations; reported serving gains include up to $\sim$3.47$\times$ time-to-first-token in late conversation rounds and $\sim$1.63$\times$ throughput on AMD CDNA4 MFMA kernels.

\paragraph{GPU decode co-design baseline}
BitDecoding~\cite{bitdecoding2026} co-designs low-bit KV formats (MXFP4/NVFP4) with CUDA and Tensor Core decode paths, making it the primary \emph{GPU-side} mixed-precision KV reference for bandwidth-bound decode.
Unlike algorithm-only PTQ papers, BitDecoding reports decode throughput, bytes moved per token, and accuracy jointly, and therefore anchors cross-platform comparisons in Table~\ref{tab:l2} even when ASIC datapaths are evaluated only architecturally.

\subsection{Hardware Mapping Patterns}
\label{sec:l2-hw}

Translating L2 algorithms to accelerators reduces to a small set of datapath choices.

\paragraph{Low-bit MAC with fused dequantization}
Ideal KV datapaths widen $Q$ (or partial sums) to FP16/FP32, fetch compressed $K/V$ tiles, apply group scales in the MAC front end, and never write full-precision caches back to HBM~\cite{liu2024kivi,sawint42026,plena2025}.
GPU implementations realize this through fused kernels (SAW-INT4, BitDecoding, UltraQuant); ASIC proposals are sparse outside PLENA's asymmetric Pathway~2/3 descriptions~\cite{plena2025}.

\paragraph{Scale and metadata management}
Per-token scales suit decode-heavy KV caches but increase metadata traffic; per-channel scales suit $K$ tensors with channel outliers; microscaling packs exponents per block for MXFP4 compatibility~\cite{bitdecoding2026,batquant2026}.
Hardware must buffer scale vectors alongside tile fetch and broadcast them to PE arrays without stalling FlashAttention-style pipelines (Section~\ref{sec:l1}).

\paragraph{Rotation and layout fusion}
QuaRot/SpinQuant rotations and SAW-INT4's BDR imply additional shuffle or butterfly network stages unless compiled into weight/K layout offline~\cite{ashkboos2024quarot,sawint42026}.
SAW-INT4 demonstrates that rotation-aware \emph{paged} KV layouts are necessary for production serving; monolithic accelerators must expose analogous address interleaving in DMA descriptors (Section~\ref{sec:l4}).

Table~\ref{tab:l2} compares representative L2 works along format, target tensor, platform, and evaluation emphasis.
As in Section~\ref{sec:l1}, absolute throughput numbers are platform-specific; we highlight whether each design primarily saves memory, bandwidth, or energy at decode.

\begin{table}[t]
\caption{Representative mixed-precision attention and KV-cache schemes (L2).
KV = key--value cache; rot.\ = rotation-based outlier handling.}
\label{tab:l2}
\centering
\scriptsize
\setlength{\tabcolsep}{2pt}
\renewcommand{\arraystretch}{1.05}
\resizebox{\columnwidth}{!}{%
\begin{tabular}{@{}lllll@{}}
\hline
\textbf{Work} & \textbf{Format} & \textbf{Target} & \textbf{Plat.} & \textbf{Focus} \\
\hline
SmoothQuant & W8A8 & Weights/act. & GPU & INT8 baseline \\
QuaRot/SpinQuant & INT4/8 & W/A/KV & GPU & Global rot. \\
KIVI & 2-bit asym. & KV & GPU & Memory \\
SAW-INT4 & INT4+BDR & KV & GPU & Serving \\
BitDecoding & MXFP4/NVFP4 & KV & GPU & Decode BW \\
BATQuant & MXFP4 BAT & W/A/KV & Algo. & MXFP align \\
UltraQuant & FP4 KV, FP8 Q & KV/query & GPU & Agent serving \\
PLENA & Asym.\ paths & KV+array & ASIC & L1--L4 stack \\
\hline
\end{tabular}%
}
\end{table}

\subsection{Accuracy--Efficiency Trade-offs}
\label{sec:l2-pareto}

Cross-study L2 evaluation should be read as a \emph{Pareto surface} over memory or HBM bytes per token, decode throughput, energy per token, and task accuracy (or perplexity), not as single-point speedups~\cite{sawint42026,batquant2026,bitdecoding2026,ultraquant2026}.
Algorithmic papers often sweep weight/activation/KV bit widths (e.g., W4A4, W4A8/KV8) on lm-eval or WikiText~\cite{ashkboos2024quarot,batquant2026,liu2024kivi}.
Serving-oriented papers add batch size, context length, and paged-cache fragmentation effects that pure accuracy sweeps omit~\cite{sawint42026,ultraquant2026,bitdecoding2026}.

Aligned with the unified metrics in Table~\ref{tab:metrics}, we recommend reporting at minimum:
(i)~effective KV width and scale granularity,
(ii)~HBM bytes read per decode step versus FP16 KV,
(iii)~downstream accuracy at matched model and context length, and
(iv)~whether prefill and decode kernels share the same numeric path.
Only then can L2 schemes be compared fairly against FP16 FlashAttention baselines (Section~\ref{sec:l1}) and eventual ASIC implementations.

\subsection{Open Problems}
\label{sec:l2-open}

Four gaps motivate continued L2--L1 co-design.
\emph{(i) ASIC datapath evidence.} Aside from PLENA~\cite{plena2025}, few monolithic accelerators publish KV dequantization fused with online softmax under fixed SRAM budgets.
\emph{(ii) Format heterogeneity.} MXFP4, INT4, and FP4 KV coexist without a common hardware denominator; block size mismatches with rotation granularity remain unresolved~\cite{batquant2026,sawint42026}.
\emph{(iii) Long-context drift.} Multi-turn agent caches accumulate quantization error over rounds~\cite{ultraquant2026}; static PTQ may be insufficient without periodic refresh or higher-precision anchor tokens.
\emph{(iv) End-to-end fusion.} GPU kernels demonstrate fused decode, but combining low-bit KV with in-array softmax and non-GEMM units (Section~\ref{sec:l3}) on one accelerator is still an open systems problem addressed in Section~\ref{sec:gaps}.

\section{Non-GEMM Specialized Hardware Units}
\label{sec:l3}

The third survey theme (L3) covers hardware for operators that sit outside dense GEMM yet gate end-to-end attention throughput: row-wise softmax (max, $\exp$, reciprocal), root-mean-square normalization (RMSNorm)~\cite{zhang2019root}, and rotary position embedding (RoPE) on $Q$ and $K$~\cite{su2024roformer}.
Sections~\ref{sec:l1}--\ref{sec:l2} showed that FlashAttention-native arrays and mixed-precision KV paths still stall when these stages are mapped to slow or off-chip vector units~\cite{wang2023cosa,desa2025,lin2025systolicattention}.
L3 literature therefore asks how to approximate or fuse non-GEMM functions with bounded error, minimal area, and scheduling that does not idle the main systolic mesh.

\subsection{Operator Roles and Integration Models}
\label{sec:l3-operators}

Within one Transformer layer, softmax executes \emph{online} alongside tiled $QK^\top$ (Section~\ref{sec:l1}); RMSNorm typically follows attention or feed-forward sub-blocks; RoPE rotates $Q/K$ before the $QK^\top$ product~\cite{su2024roformer,zhang2019root}.
Hardware integration follows three patterns seen across L1--L3:
\emph{(i) dedicated side units}---COSA+ places a softmax engine between dual systolic arrays~\cite{wang2023cosa};
\emph{(ii) in-array fusion}---SystolicAttention embeds row-max/row-sum and split-integer $\exp$ inside the MAC mesh~\cite{lin2025systolicattention,koca2023exp};
\emph{(iii) auxiliary vector or immediate-processing paths}---DESA routes non-GEMM work through an IPU and reconfigurable vector logic~\cite{desa2025}.
The design tension is identical in each case: specialization saves cycles but adds control complexity and risks pipeline bubbles unless compile-time schedulers overlap DMA, GEMM, and non-GEMM stages (Section~\ref{sec:l4}).

\subsection{Softmax and Exponential Approximation}
\label{sec:l3-softmax}

Because FlashAttention subtracts the row maximum before exponentiation~\eqref{eq:online-max}, all $\exp$ inputs are $\leq 0$, narrowing the approximation range for hardware~\cite{dao2022flashattention,lin2025systolicattention}.

\paragraph{Lookup and piecewise methods}
Dedicated softmax accelerators often combine max-reduction trees with ROM or bipartite lookup tables (LUTs) for $\exp$ and dividers for normalization~\cite{softmaxrmsnorm2026,guaranteednorm2026,mive2026,ham2021a3}.
An FPGA-oriented design pairs SafeSoftmax with bipartite LUT $\exp$ and reports pipelined latency plus DSP/LUT cost~\cite{softmaxrmsnorm2026}.
Guaranteed-normalization softmax uses radix-2 $\exp$ LUTs with fixed-point division to preserve \emph{absolute} scale---not just ordering---which matters for score-sensitive tasks where order-only approximations fail~\cite{guaranteednorm2026}; authors report ASIC areas of $\sim$942\,$\mu$m$^2$ for softmax and $\sim$1199\,$\mu$m$^2$ for LayerNorm at $N{+}1$ cycle latency, with $\sim$11$\times$/14$\times$ smaller area than reference blocks.

\paragraph{Base-2 and log-domain variants}
Softermax replaces $e^x$ with $2^x$ to implement softmax via shifts and adders~\cite{sun2022softmax}, yielding $\sim$0.25$\times$ area versus a floating $\exp$ unit on FPGA at moderate accuracy cost.
SOLE's E2Softmax quantizes $\exp$ outputs to 4-bit in a $\log_2$ domain using a Log2Exp unit built from shifters and adders, avoiding full multipliers~\cite{sole2025}.

\paragraph{In-array and MAC-reuse approaches}
SystolicAttention splits each non-positive input into integer and fractional parts and linearly interpolates $\exp$ while reusing systolic MAC hardware~\cite{lin2025systolicattention,koca2023exp}.
StreamAttention adopts Chebyshev-polynomial $\exp$ inside a streaming single-array schedule~\cite{streamattention2025}.
These approaches minimize off-array traffic but tie numeric error to tile scheduling; errors compound through online softmax accumulation (Section~\ref{sec:l2}).

\subsection{RMSNorm and Reciprocal Square Root}
\label{sec:l3-rmsnorm}

RMSNorm computes $x_i \cdot \mathrm{rsqrt}(\operatorname{mean}(x^2)+\epsilon)$ and requires sum-of-squares reduction followed by inverse square root~\cite{zhang2019root}.
Hardware implementations mirror softmax trends: reduction trees plus approximate $\mathrm{rsqrt}$.
The FPGA design in~\cite{softmaxrmsnorm2026} uses LOD-LUT-MUL for $\mathrm{rsqrt}$;~\cite{guaranteednorm2026} applies Newton iterations for LayerNorm-style normalization with bounded cycle count.
SOLE's AILayerNorm shares low-precision statistics with E2Softmax, reporting $\sim$3.04$\times$/$\sim$3.86$\times$ energy-efficiency and $\sim$2.82$\times$/$\sim$3.32$\times$ area-efficiency improvements over separate softmax/LayerNorm accelerators on FPGA/ASIC~\cite{sole2025}.
MIVE unifies softmax, LayerNorm, and RMSNorm in one minimalist integer vector engine with shared ROM LUTs for $\exp$ and reciprocal, improving area efficiency relative to three standalone units~\cite{mive2026}.

\subsection{RoPE and Remaining Gaps}
\label{sec:l3-rope}

RoPE applies position-dependent rotations to $Q$ and $K$ using $\cos/\sin$ lookups and complex-style pairwise multiplies~\cite{su2024roformer}.
Unlike softmax and RMSNorm, few L3 papers treat RoPE as a first-class accelerator block; most full-stack proposals assume FP16/FP32 LUTs fused into projection datapaths without independent Pareto reporting~\cite{plena2025,wang2023cosa,desa2025}.
% REVIEW: dedicated RoPE ASIC literature remains sparse as of 2026 survey cutoff
This remains a documented blind spot: long-context serving shifts bottleneck to KV traffic (Section~\ref{sec:l2}), but RoPE still executes once per layer on $Q/K$ and can become latency-visible when GEMM is well optimized.

\subsection{Unified Datapaths and Comparison}
\label{sec:l3-unified}

Recent work converges on \emph{shared} rather than per-operator accelerators.
MIVE programmable integer datapaths cover softmax, LayerNorm, and RMSNorm with one ROM-heavy engine~\cite{mive2026}.
SOLE co-designs E2Softmax and AILayerNorm to keep intermediate statistics on chip rather than spilling memory-bound tensors between units~\cite{sole2025}.
At the opposite extreme, in-array softmax fusion (SystolicAttention, StreamAttention) reduces unit count to zero at the cost of tighter coupling to GEMM tiling~\cite{lin2025systolicattention,streamattention2025}.

Table~\ref{tab:l3} summarizes representative L3 approaches.
Evaluation spans FPGA demonstrations, ASIC synthesis, and RTL modules co-reported with L1 arrays; cosine similarity or task accuracy versus FP32 should accompany any approximate datapath claim (Table~\ref{tab:metrics}).

\begin{table}[t]
\caption{Representative non-GEMM hardware (L3).
PWL = piecewise linear; int.\ = in-array integration.}
\label{tab:l3}
\centering
\scriptsize
\setlength{\tabcolsep}{2pt}
\renewcommand{\arraystretch}{1.05}
\resizebox{\columnwidth}{!}{%
\begin{tabular}{@{}lllll@{}}
\hline
\textbf{Work} & \textbf{Operators} & \textbf{Approx.} & \textbf{Integration} & \textbf{Eval.} \\
\hline
Softermax & Softmax & Base-$2$ $\exp$ & Side / FPGA & FPGA \\
Koca et al. & $\exp$ & Split-int.{+}PWL & MAC reuse & RTL \\
SystolicAttn & Softmax & In-array $\exp$ & Systolic fuse & RTL \\
Softmax/RMSNorm & Soft.+RMSNorm & LUT + LOD rsqrt & Standalone & FPGA \\
Guar.\ norm & Soft.+LayerNorm & Radix-2 LUT & Standalone & ASIC \\
SOLE & Soft.+LayerNorm & E2Softmax/log & Unified unit & FPGA/ASIC \\
MIVE & Soft.+LN/RMSNorm & Integer LUT & Unified engine & ASIC \\
COSA+/DESA & Softmax (+norm) & Mixed & Side IPU/unit & Arch. \\
\hline
\end{tabular}%
}
\end{table}

\subsection{Open Problems}
\label{sec:l3-open}

Three L3 gaps connect directly to monolithic accelerator goals.
\emph{(i) Error composition.} Approximate softmax, low-bit KV dequantization (Section~\ref{sec:l2}), and RMSNorm rounding interact across layers; few studies publish joint sensitivity analysis~\cite{guaranteednorm2026,sole2025,liu2024kivi}.
\emph{(ii) Pipeline orchestration.} Unified engines (MIVE, SOLE) save area but require explicit latency matching with systolic arrays; in-array fusion avoids handoff yet limits reuse for feed-forward norms~\cite{mive2026,lin2025systolicattention,desa2025}.
\emph{(iii) RoPE and beyond.} Dedicated RoPE hardware, fused Q/K projection--rotation datapaths, and standardized accuracy metrics for positional encoding remain largely unaddressed~\cite{su2024roformer,plena2025}.
Closing these gaps alongside L1 dataflow and L2 precision choices is necessary for the integrated co-design agenda in Section~\ref{sec:gaps}.

\section{Compilation, Mapping, and Software-Hardware Co-Design}
\label{sec:l4}

The fourth survey theme (L4) covers compile-time and toolchain support needed to realize the L1--L3 datapaths on monolithic accelerators: FlashAttention tiling, on-chip buffer allocation, direct memory access (DMA) scheduling, mixed-precision configuration, and power--performance--area (PPA) assessment.
GPU stacks implement these concerns implicitly in kernel schedulers~\cite{dao2022flashattention,shah2024flashattention3,sawint42026}; ASIC literature increasingly exposes them through custom instruction-set architectures (ISAs), design-space exploration (DSE), and cycle-accurate simulators~\cite{plena2025,flatattention2026,klhufek2025transinfersim}.
Without L4 automation, fixed hardware for online softmax (Section~\ref{sec:l1}), KV dequantization (Section~\ref{sec:l2}), and non-GEMM fusion (Section~\ref{sec:l3}) cannot adapt across prefill/decode phases or context lengths.

\subsection{Tiling, Buffers, and Memory Scheduling}
\label{sec:l4-tiling}

FlashAttention-style execution requires joint selection of tile sizes along sequence length and head dimension, SRAM capacity per buffer ($Q$, $K$, $V$, partial $O$, running softmax statistics), and double-buffering depth to overlap HBM fetch with array compute~\cite{dao2022flashattention,lin2025systolicattention,plena2025}.
Classical DNN mappers such as Timeloop formulate tiling as constrained optimization over loop nests and memory hierarchies~\cite{parashar2019timeloop,chen2016eyeriss}; attention adds \emph{phase-dependent} constraints because prefill uses large parallel $Q$ blocks whereas decode streams a single query row against a growing KV cache (Sections~\ref{sec:l1-prefill-decode},~\ref{sec:l2}).

\paragraph{Phase-adaptive mapping}
PLENA's compiler rebalances head-level decomposition and asymmetric precision pathways between prefill and decode on a flattened array~\cite{plena2025}.
FlatAttention co-optimizes attention loop nests with on-chip \emph{fabric collectives}, choosing multicast/reduction patterns across tiles to cut cross-die HBM traffic rather than optimizing tiles in isolation~\cite{flatattention2026}.
Both illustrate that attention mappers must co-design \emph{dataflow and interconnect}, not only PE array dimensions.

\paragraph{DMA--compute overlap}
Effective L4 schedules interleave DMA transfers for $K/V$ tiles (including compressed KV metadata from Section~\ref{sec:l2}) with systolic execution and non-GEMM side units~\cite{desa2025,wang2023cosa,plena2025}.
Paged-KV serving layouts on GPUs~\cite{sawint42026,bitdecoding2026} imply scatter-gather DMA descriptors on ASIC; compile-time lowering must emit physical addresses and scale pointers consistent with rotation or microscaling block boundaries~\cite{ashkboos2024quarot,batquant2026}.

\subsection{Custom ISA and Kernel Lowering}
\label{sec:l4-isa}

Full-stack accelerator proposals bind hardware to a programmable interface.
PLENA provides a custom ISA, compiler, transaction-level simulator, and RTL flow spanning FlashAttention-native ops, mixed-precision MAC modes, and long-context agentic benchmarks~\cite{plena2025}; it is the primary L4 reference for \emph{integrated} L1--L2 mapping on ASIC.
At the GPU end of the spectrum, TileLang offers a Python tile DSL with explicit memory-hierarchy buffers for kernel authors who need finer control than vendor libraries provide~\cite{tilelang2025}; it supports rapid FlashAttention-variant prototyping but does not by itself solve ASIC DMA scheduling.

The recurring lowering pipeline across L4 systems is:
\emph{(i)} graph-level partitioning into attention, norm, and projection subgraphs;
\emph{(ii)} loop tiling and buffer allocation subject to SRAM caps;
\emph{(iii)} instruction or kernel emission with fused dequant/softmax flags; and
\emph{(iv)} phase-specific replanning when sequence length or batch size changes between prefill and decode serving modes~\cite{plena2025,flatattention2026,sawint42026}.

\subsection{Evaluation Toolchains}
\label{sec:l4-tools}

Surveyed accelerators are evaluated at three calibrated layers (Table~\ref{tab:metrics}): algorithmic accuracy, architectural PPA, and RTL synthesis.
L4 tools anchor the architectural layer and interface to the others.

\paragraph{Analytical mapping and energy}
Timeloop searches mappings with analytic latency/energy models~\cite{parashar2019timeloop}; Accelergy supplies component-level energy-per-action tables for SRAM, MAC, and controller accesses~\cite{wu2019accelergy}.
Together they remain the default backbone for comparing tiling choices and buffer hierarchies when cycle simulators are unavailable.
Eyeriss established row-stationary dataflow co-design methodology that informs modern systolic mappers~\cite{chen2016eyeriss}.

\paragraph{Cycle-accurate systolic simulation}
SCALE-Sim v3 extends cycle-accurate systolic modeling with multi-core execution, sparsity hooks, Ramulator DRAM, and Accelergy integration for end-to-end bandwidth and energy traces~\cite{scalesimv3_2025}.
TransInferSim targets Transformer inference specifically, modeling layer-level execution on systolic arrays with cache hierarchies, exporting execution plans, and integrating Accelergy for area/energy estimates~\cite{klhufek2025transinfersim}.
Both tools help validate whether a proposed FlashAttention tile schedule sustains PE utilization reported in L1 papers, but require user-supplied models for online softmax and non-GEMM units (Section~\ref{sec:l3}).

\paragraph{Algorithmic and RTL anchors}
Algorithmic evaluation uses PyTorch reference kernels, lm-eval accuracy, and cosine-similarity checks against FP32 golden models~\cite{ashkboos2024quarot,liu2024kivi,guaranteednorm2026}.
RTL-level modules---for example SystolicAttention's 16\,nm synthesis~\cite{lin2025systolicattention}---ground cycle counts and area overhead for fused softmax datapaths.
Cross-layer calibration demands that the \emph{same} tile dimensions and precision modes flow from algorithm scripts through simulators to RTL testbenches; mismatches here explain contradictory utilization claims across L1--L3 studies.

Table~\ref{tab:l4} summarizes representative L4 artifacts.
Open-source availability (Timeloop, Accelergy, SCALE-Sim v3, TransInferSim) lowers the barrier to reproducible cross-study comparison relative to closed GPU serving stacks.

\begin{table}[t]
\caption{Representative compilation, mapping, and evaluation tools (L4).
TL = Timeloop; Acc = Accelergy; cycle-acc.\ = cycle-accurate.}
\label{tab:l4}
\centering
\scriptsize
\setlength{\tabcolsep}{2pt}
\renewcommand{\arraystretch}{1.05}
\resizebox{\columnwidth}{!}{%
\begin{tabular}{@{}lllll@{}}
\hline
\textbf{Work} & \textbf{Role} & \textbf{Target} & \textbf{Layer} & \textbf{Open} \\
\hline
Timeloop & Mapping search & Generic array & Arch. & Yes \\
Accelergy & Energy tables & Components & Arch. & Yes \\
SCALE-Sim v3 & Systolic sim & Multi-core array & Cycle-acc. & Yes \\
TransInferSim & Transformer sim & Systolic+cache & Cycle-acc. & Yes \\
TileLang & Tile DSL & GPU/NPU kernels & Kernel & TBD \\
PLENA stack & ISA+compiler & PLENA ASIC & Arch.+RTL & Planned \\
FlatAttention & Fabric+tile DSE & Tile accelerator & Arch. & No \\
Eyeriss & Row-stationary & Systolic chip & Silicon & No \\
\hline
\end{tabular}%
}
\end{table}

\subsection{Open Problems}
\label{sec:l4-open}

Four L4 gaps persist despite mature DNN mapping tools.
\emph{(i) Attention-native cost models.} Timeloop/SCALE-Sim default to GEMM-centric loop nests; online softmax, KV dequantization, and non-GEMM latency are usually overlaid manually~\cite{parashar2019timeloop,scalesimv3_2025,lin2025systolicattention,mive2026}.
\emph{(ii) Phase-aware autotuning.} Few open tools jointly optimize prefill and decode schedules under a single SRAM budget~\cite{plena2025,dao2023flashdecoding,flatattention2026}.
\emph{(iii) Precision-aware mapping.} Mixed-precision KV (Section~\ref{sec:l2}) changes tile footprints and DMA width; public mappers seldom co-search bit width and loop order~\cite{batquant2026,sawint42026,plena2025}.
\emph{(iv) Reproducible end-to-end flows.} GPU serving papers rarely release mappers; ASIC papers seldom link algorithmic, simulator, and RTL artifacts in one repository~\cite{plena2025,klhufek2025transinfersim,sawint42026}.
Unified L4 infrastructure that spans Sections~\ref{sec:l1}--\ref{sec:l3} remains a prerequisite for the integrated research agenda in Section~\ref{sec:gaps}.

\section{Adjacent Directions and Outlook}
\label{sec:adjacent}

Sections~\ref{sec:l1}--\ref{sec:l4} focused on \emph{monolithic} Transformer inference accelerators with explicit attention datapath co-design.
Adjacent lines of work attack similar bottlenecks---long-context KV traffic and attention latency---through sparsity, dynamic token pruning, or memory-centric scale-out.
We summarize them here to clarify survey boundaries (Section~\ref{sec:intro-scope}): they inform outlook and baseline comparisons but are not primary taxonomy themes.

\subsection{Sparse and Dynamic Attention}
\label{sec:adjacent-sparse}

Sparse attention reduces effective sequence length by skipping tokens, heads, or blocks whose contributions fall below a threshold.
SpAtten~\cite{wang2021spatten} prunes tokens and heads in a cascade and builds dedicated sparse-attention hardware for energy savings.
Sanger~\cite{liu2021sanger} co-designs a reconfigurable systolic array with score-stationary dataflow and quantization-aware sparsity prediction, reporting substantial speedups over dense attention accelerators.
Salca~\cite{salca2026} targets long-context \emph{decode} with dual-compression dynamic sparse attention, combining ultra-low precision and feature sparsity in a pipelined ASIC; authors report large speedup and energy gains versus GPU baselines at the cost of additional control for sparsity pattern generation and load balance.

These designs complement dense FlashAttention-native arrays (Section~\ref{sec:l1}) when attention maps are structurally sparse or when approximate top-$k$ selection preserves task accuracy.
They introduce new compiler and scheduling problems---irregular memory access, load imbalance across PEs, and interaction with low-bit KV formats (Section~\ref{sec:l2})---that dense monolithic mappers do not yet model~\cite{liu2021sanger,salca2026,parashar2019timeloop}.
Block-sparse GPU kernels such as the block-sparse variant of FlashAttention~\cite{dao2022flashattention} provide algorithmic references but assume programmer-specified sparsity patterns rather than hardware-driven pruning.

\subsection{Memory-Centric and Heterogeneous Scale-Out}
\label{sec:adjacent-memory}

When context length or batch size exceeds single-die HBM capacity and bandwidth, architects move attention state closer to memory or across chiplets.
AMMA~\cite{amma2026} proposes a multi-chiplet memory-centric organization with HBM processing-near-memory cubes for million-token context serving, reporting large latency and energy improvements over H100-class GPUs at the expense of multi-die packaging complexity.
LoL-PIM~\cite{lolpim2024} pipelines long-context decode across DRAM-PIM nodes with an MLIR-based compiler when KV caches exceed local PIM capacity.
NeuPIM~\cite{heo2024nuepim} couples NPUs with PIM for batched LLM inference, offloading bandwidth-heavy KV movement to memory-side compute.

Such systems address the same memory-wall motivation as PLENA and BitDecoding but at \emph{system} granularity.
They are out of scope for the monolithic L1--L4 taxonomy yet set upper-bound targets for bandwidth reduction that single-chip designs must approach through tiling, mixed precision, and on-chip fusion rather than through package-scale disaggregation~\cite{plena2025,bitdecoding2026,amma2026}.

\subsection{Boundary and Outlook}
\label{sec:adjacent-outlook}

We treat distributed training, mixture-of-experts (MoE) routing, compute-in-memory (CIM) research without Transformer attention co-design, and multi-chiplet scale-out as adjacent rather than core topics.
Within the monolithic scope, the dominant trend remains \emph{dense} FlashAttention-native execution with progressively aggressive KV compression and fused non-GEMM datapaths.
Sparse and memory-centric directions are most relevant as \emph{fallbacks} when dense arrays exhaust SRAM or HBM budgets, or as comparators for energy-per-token at extreme context.
Closing the loop between dense single-chip co-design (Sections~\ref{sec:l1}--\ref{sec:l4}) and these adjacent stacks---for example, exporting the same KV layout to PIM nodes or gating sparsity from compile-time profiles---is an open systems integration question beyond current surveyed monolithic proposals.

\section{Research Gaps and Positioning}
\label{sec:gaps}

Cross-review of Sections~\ref{sec:l1}--\ref{sec:l4} confirms the survey storyline: LLM inference is increasingly memory-bound on attention and KV traffic~\cite{dao2023flashdecoding,bitdecoding2026,sawint42026}; L1--L4 each mitigate one facet; yet most publications optimize one or two facets in isolation on mismatched platforms and benchmarks~\cite{lin2025systolicattention,liu2024kivi,mive2026,plena2025,sawint42026}.
Table~\ref{tab:metrics} unifies reporting dimensions, but phase-separated prefill/decode data, open mapper artifacts, and joint accuracy--PPA curves remain rare.
This section synthesizes research gaps and relates them to the integrated monolithic accelerator program outlined in the companion research plan.

\subsection{Cross-Cutting Research Gaps}
\label{sec:gaps-synthesis}

\paragraph{Incomplete single-chip co-design}
GPU kernels (FlashAttention, BitDecoding, SAW-INT4) optimize IO-aware dataflow and KV compression without custom silicon~\cite{dao2022flashattention,bitdecoding2026,sawint42026}.
SystolicAttention/StreamAttention (single-array) and COSA+ (dual-array) optimize softmax mapping but largely assume FP16 internal precision and omit systematic KV quantization studies~\cite{lin2025systolicattention,streamattention2025,wang2023cosa}.
Non-GEMM units (MIVE, SOLE) and mixed-precision algorithms (KIVI, BATQuant) are rarely co-evaluated with FlashAttention tiling and compile-time schedulers on the same platform~\cite{mive2026,sole2025,liu2024kivi,batquant2026,plena2025}.
PLENA is the closest end-to-end exception but remains the only widely cited monolithic stack spanning L1--L4~\cite{plena2025}.

\paragraph{Phase and precision coupling}
Prefill favors wide $QK^\top$ tiles; decode favors KV bandwidth and skinny matmul schedules (Sections~\ref{sec:l1-prefill-decode},~\ref{sec:l2}).
Few mappers co-optimize both under one SRAM cap or co-search bit width, tile order, and DMA depth~\cite{plena2025,flatattention2026,dao2023flashdecoding,batquant2026}.
Online softmax invariants (Section~\ref{sec:l1-online}) constrain how aggressively KV may be compressed without destabilizing accumulation~\cite{liu2024kivi,guaranteednorm2026}.

\paragraph{Evaluation fragmentation}
Architectural simulators default to GEMM-centric loop nests~\cite{parashar2019timeloop,scalesimv3_2025,klhufek2025transinfersim}; RTL modules validate subsystems in isolation~\cite{lin2025systolicattention,mive2026}.
Without three-layer calibration---algorithmic golden models, cycle-accurate mapping, and RTL PPA---reported utilization and energy-per-token figures are difficult to compare across L1--L3 (Section~\ref{sec:l4-tools}).

\subsection{Open Problem Checklist}
\label{sec:gaps-open}

The following open problems recur across surveyed themes:
\begin{enumerate}
  \item \textbf{Numerical stability under mixed precision:} joint bounds on KV quantization, softmax approximation, and RMSNorm error across long contexts~\cite{guaranteednorm2026,sole2025,batquant2026,ultraquant2026}.
  \item \textbf{Non-GEMM error composition:} end-to-end sensitivity when approximate $\exp/\mathrm{rsqrt}$ units (Section~\ref{sec:l3-softmax},~\ref{sec:l3-rmsnorm}) interact with low-bit KV dequantization~\cite{lin2025systolicattention,liu2024kivi,mive2026}.
  \item \textbf{Dual-mode scheduling:} unified compile-time plans for prefill and decode with replanning at serving time~\cite{plena2025,streamattention2025,dao2023flashdecoding}.
  \item \textbf{RoPE and projection fusion:} lack of Pareto-validated hardware for positional encoding (Section~\ref{sec:l3-rope})~\cite{su2024roformer,plena2025}.
  \item \textbf{Attention-native cost models:} mapper support for online softmax, metadata scales, and fabric/collective effects~\cite{flatattention2026,parashar2019timeloop,scalesimv3_2025}.
  \item \textbf{Reproducible artifact chains:} open releases linking algorithm scripts, simulators, and RTL for monolithic attention accelerators~\cite{plena2025,klhufek2025transinfersim,sawint42026}.
\end{enumerate}

\subsection{Positioning Relative to Prior Work}
\label{sec:gaps-positioning}

The gaps above motivate an integrated monolithic accelerator agenda---not incremental improvement within a single L-theme---with three differentiation vectors aligned to prior art:

\paragraph{Architecture (L1)}
Target a FlashAttention-\emph{native} systolic or flattened array with \emph{in-place} online softmax (row-max/row-sum and fused $\exp$), eliminating full $N \times N$ score materialization and reducing dependence on external vector softmax units~\cite{lin2025systolicattention,wang2023cosa,desa2025}.
Unlike COSA+'s dual-array partition or GPU kernels tied to shared memory, the goal is single-chip utilization across prefill \emph{and} decode through phase-aware mapping~\cite{plena2025,streamattention2025}.

\paragraph{Numerics and datapath (L2--L3)}
Combine mixed-precision KV storage (INT4/FP8/MXFP4 class formats) with higher-precision matmul/softmax accumulators, rotation- or block-aware outlier suppression, and on-chip dequantization fused at the MAC boundary~\cite{sawint42026,batquant2026,liu2024kivi,ashkboos2024quarot}.
Pair this with unified or in-array non-GEMM approximations whose error budgets are validated jointly rather than per operator~\cite{mive2026,sole2025,guaranteednorm2026,koca2023exp}.

\paragraph{Software--hardware co-design (L4)}
Provide automatic tiling, buffer allocation, mixed-precision configuration, and GEMM/non-GEMM pipeline scheduling with a custom ISA or intermediate representation, using Timeloop/Accelergy, SCALE-Sim v3, and TransInferSim for architectural exploration and RTL for calibration~\cite{parashar2019timeloop,wu2019accelergy,scalesimv3_2025,klhufek2025transinfersim,plena2025}.

\paragraph{Quantifiable targets}
% TODO: calibrate against FlashAttention/FlashDecoding and PLENA baselines after phase-1 modeling
Relative to GPU FlashAttention/FlashDecoding and documented monolithic references, quantifiable goals for such an integrated design include improved PE utilization, reduced HBM bytes per decode token, and lower energy per token at 32K--128K context, with task accuracy or perplexity degradation bounded by a Pareto envelope comparable to SAW-INT4/KIVI-class KV schemes~\cite{sawint42026,liu2024kivi,dao2023flashdecoding}.
Absolute targets remain \emph{to be calibrated} once baseline models and workloads are fixed (Table~\ref{tab:metrics}).

Adjacent sparse and memory-centric systems (Section~\ref{sec:adjacent}) serve as upper-bound comparators and scope guards rather than primary implementation paths.
Success for the monolithic agenda is demonstrated by closing the L1--L4 gaps \emph{on one platform} with reproducible evaluation artifacts, not by matching multi-chiplet or PIM peak bandwidth in isolation~\cite{amma2026,salca2026,plena2025}.

\section{Conclusion}
\label{sec:conclusion}

Transformer inference has shifted from raw compute throughput to attention and key-value (KV) cache data movement as the dominant bottleneck in long-context and decode-heavy workloads.
This survey reviewed monolithic accelerator research through four complementary themes: FlashAttention-native dataflows and array organizations (Section~\ref{sec:l1}), mixed-precision attention datapaths and KV-cache quantization (Section~\ref{sec:l2}), specialized hardware for non-general-matrix-multiply (non-GEMM) operators (Section~\ref{sec:l3}), and compilation, mapping, and evaluation toolchains (Section~\ref{sec:l4}).
Across these themes, IO-aware online softmax, low-width KV storage with higher-precision accumulation, approximate softmax and normalization units, and phase-aware mappers emerge as recurring design patterns, while GPU FlashAttention/FlashDecoding kernels remain the primary algorithmic references for correctness and traffic complexity.

Our synthesis highlights a persistent fragmentation of the design space.
Most studies optimize one or two concerns---for example, GPU decode bandwidth, single-array softmax fusion, KV post-training quantization, or standalone normalization accelerators---on incomparable platforms and benchmarks.
Unified comparison along shared metrics (Table~\ref{tab:metrics}) and separation of prefill versus decode results are necessary but still inconsistently applied.
Adjacent sparse-attention and memory-centric scale-out directions (Section~\ref{sec:adjacent}) address related memory-wall pressures at system granularity but lie outside the monolithic taxonomy developed here.

Section~\ref{sec:gaps} argued that the central open problem is \emph{integrated} co-design on a single chip: FlashAttention-native tiling, mixed-precision KV datapaths, fused non-GEMM stages, and compile-time scheduling validated jointly through algorithmic, architectural, and RTL evaluation.
Until such end-to-end artifact chains become common, reported energy-per-token and utilization gains will remain difficult to interpret across the L1--L4 literature.
Closing this gap is the logical next step for energy-efficient attention acceleration research within the monolithic scope defined at the outset of this survey.

\bibliographystyle{IEEEtran}
\begingroup
\raggedright
\bibliography{references}
\endgroup

\end{document}
